\documentclass{plantillaPPT}
\titulo{13}{Series de potencia y Taylor}
\begin{document}
\primeraDiapo

\begin{frame}{Problema 1}
\framesubtitle{Práctica 13}

    1. Sabiendo que $\displaystyle \forall x\in\mathbb{R}, e^x = \sum_{n=1}^\infty \frac{x^n}{n!}$ exprese $f(x) = xe^x$ como una serie y muestre que
    \begin{equation*}
        \sum_{n=1}^\infty \frac{(-1)^n}{(n+2)n!} = \frac{1}{2}-\frac{2}{e}
    \end{equation*}
\end{frame}

\begin{frame}{Series de Taylor}
\framesubtitle{Aproximación de funciones}

% --- Definición ---
\begin{definicion}[Serie de Taylor]
    Sea \(f\) una función infinitamente diferenciable en un intervalo \((a, b)\) que contiene al punto \(c\).
    La \textbf{Serie de Taylor} de \(f\) centrada en \(c\) es la serie de potencias dada por:
\end{definicion}

% --- Expresión de la Serie ---
\begin{equation*}
f(x) = \sum_{n=0}^{\infty} \frac{f^{(n)}(c)}{n!}(x-c)^n
\end{equation*}

% --- Parte Extra/Comentario ---
\begin{itemize}
    \item Esta serie representa a la función \(f(x)\) en el \textbf{radio de convergencia} de la serie.
    \item Si el centro es \(c=0\), la serie se denomina \textbf{Serie de Maclaurin}.
\end{itemize}

\end{frame}

\begin{frame}{Problema 2}
    \framesubtitle{Práctica 13}
    
    2. Determine la serie de Taylor centrada en $c$, para la función $f$ indicada.\\[15pt]
    $(a)\quad f(x)=\ln(4x+3),\; c=0$.\\
    $(b)\quad f(x) = \frac{1}{x^4},\; c=-3$. \\
    $(c)\quad f(x) = \arctan(x),\; c=0$.
\end{frame}

\begin{frame}{Problema 3}
    \framesubtitle{Práctica 13}
    3. La integral
    \begin{equation*}
        I = \int_0^1 (e^{x^2} - e^{-x^2})\,dx
    \end{equation*}
    no se puede calcular por métodos elementales. Nuestro objetivo es aproximar su valor utilizando series de Maclaurin.\\[10pt]
    $(a)$ Obtener la serie de Maclaurin de la función $g(x) = e^{x^2}-e^{-x^2}$ si se sabe que:
    \begin{equation*}
        \forall x\in\mathbb{R}: \frac{1}{2}(e^x-e^{-x}) = \sum_{n=0}^\infty \frac{x^{2n+1}}{(2n+1)!}
    \end{equation*}
    
\end{frame}
\begin{frame}{Problema 3}
    \framesubtitle{Práctica 13}
    $(b)$ Determinar la serie numérica que represente a la integral \\[5pt] $\displaystyle \int_0^x (e^{t^2}-e^{-t^2})\,dt$.\\[5pt]
    $(c)$ Usar los primeros tres términos de la serie obtenida en $(b)$ para aproximar el valor de $I$.
\end{frame}
\end{document}


